\subsection{Производные высших порядков. Правило Лейбница.}

\subsubsection*{Определение производных высших порядков}

Пусть функция $f(x)$ дифференцируема в некотором промежутке.
Производная второго порядка определяется как производная от первой производной:
\[ f''(x) = (f'(x))' \]
Аналогично, производная $n$-го порядка определяется рекуррентно:
\[ f^{(n)}(x) = \left( f^{(n-1)}(x) \right)' \]
(при условии, что $f^{(n-1)}(x)$ определена и дифференцируема).

\textbf{Обозначения:}
\[ f^{(n)}(x) \quad \text{или} \quad \frac{d^n f}{dx^n} \]
При этом полагают $f^{(0)}(x) = f(x)$.

\subsubsection*{Формула Лейбница}

\begin{tcolorbox}[title=Теорема (Правило Лейбница)]
    Пусть функции $u(x)$ и $v(x)$ имеют производные до порядка $n$ включительно на интервале $(a, b)$.
    Тогда их произведение $u(x)v(x)$ также $n$ раз дифференцируемо, и справедлива формула:
    \[ (uv)^{(n)} = \sum_{k=0}^n C_n^k u^{(k)} v^{(n-k)} \]
    где $C_n^k = \frac{n!}{k!(n-k)!}$ — биномиальные коэффициенты, $u^{(k)}$ — производная порядка $k$.
\end{tcolorbox}

\textbf{Доказательство (методом математической индукции):}

\textbf{1. База индукции ($n=1$):}
Формула дает:
\[ (uv)^{(1)} = C_1^0 u^{(0)}v^{(1)} + C_1^1 u^{(1)}v^{(0)} = 1 \cdot u \cdot v' + 1 \cdot u' \cdot v \]
Это совпадает с обычным правилом дифференцирования произведения: $(uv)' = u'v + uv'$. База верна.

\textbf{2. Индукционный переход ($m \to m+1$):}
Предположим, что формула верна для $n=m$:
\[ (uv)^{(m)} = \sum_{k=0}^m C_m^k u^{(k)} v^{(m-k)} \]
Найдем производную порядка $m+1$, продифференцировав это равенство:
\[ (uv)^{(m+1)} = \left( (uv)^{(m)} \right)' = \left( \sum_{k=0}^m C_m^k u^{(k)} v^{(m-k)} \right)' \]
В силу линейности производной и правила производной произведения:
\[ = \sum_{k=0}^m C_m^k \left( (u^{(k)})' v^{(m-k)} + u^{(k)} (v^{(m-k)})' \right) = \]
\[ = \sum_{k=0}^m C_m^k \left( u^{(k+1)} v^{(m-k)} + u^{(k)} v^{(m+1-k)} \right) \]

Разобьем эту сумму на две:
\[ S = \underbrace{\sum_{k=0}^m C_m^k u^{(k+1)} v^{(m-k)}}_{S_1} + \underbrace{\sum_{k=0}^m C_m^k u^{(k)} v^{(m+1-k)}}_{S_2} \]

В первой сумме $S_1$ сделаем замену индекса суммирования: пусть $j = k+1$.
Тогда при $k=0 \Rightarrow j=1$, при $k=m \Rightarrow j=m+1$.
\[ S_1 = \sum_{j=1}^{m+1} C_m^{j-1} u^{(j)} v^{(m+1-j)} \]
Вторую сумму $S_2$ оставим без изменений (обозначив индекс буквой $j$ для удобства):
\[ S_2 = \sum_{j=0}^m C_m^j u^{(j)} v^{(m+1-j)} \]

Теперь выделим слагаемые при $j=0$ и $j=m+1$:
\begin{itemize}
    \item Из $S_2$ берем слагаемое при $j=0$: $C_m^0 u^{(0)} v^{(m+1)} = 1 \cdot u v^{(m+1)}$.
    \item Из $S_1$ берем слагаемое при $j=m+1$: $C_m^m u^{(m+1)} v^{(0)} = 1 \cdot u^{(m+1)} v$.
    \item Для $j$ от $1$ до $m$ объединяем слагаемые:
    \[ u^{(j)} v^{(m+1-j)} \left( C_m^{j-1} + C_m^j \right) \]
\end{itemize}

Используем свойство биномиальных коэффициентов (тождество Паскаля):
\[ C_m^{j-1} + C_m^j = \frac{m!}{(j-1)!(m-j+1)!} + \frac{m!}{j!(m-j)!} = \frac{m!(j + m - j + 1)}{j!(m+1-j)!} = \frac{(m+1)!}{j!(m+1-j)!} = C_{m+1}^j \]

Таким образом:
\[ (uv)^{(m+1)} = \underbrace{C_{m+1}^0}_{1} u^{(0)}v^{(m+1)} + \sum_{j=1}^m C_{m+1}^j u^{(j)} v^{(m+1-j)} + \underbrace{C_{m+1}^{m+1}}_{1} u^{(m+1)}v^{(0)} \]
Сворачиваем обратно в одну сумму:
\[ (uv)^{(m+1)} = \sum_{k=0}^{m+1} C_{m+1}^k u^{(k)} v^{(m+1-k)} \]
Что и требовалось доказать.


